% Sol, step-15 manuscript contribution. New prose; locked outline unchanged.
% Requires amsmath and booktabs. Citation keys belong to phase-3 references.bib.
% Sections may be integrated separately. Numbers are finite-cohort descriptions.

\section{Study design and economic measurement}
\label{sec:sol-economic-methods}

\subsection{The cohort and the object of evaluation}

We study twenty completed four-player games, comprising eighty player-games, 50,918 events, 6,989 decisions and 8,679 response attempts. Five seeds each have four cyclic seat rotations. Sol, Gemini, Opus and Grok each occupy seats 1--4 five times. The evaluated participant is the configured agent, including its model, inference settings, observations, memory and response-resolution policy. This complete-system interpretation follows existing arguments for agent evaluation and is especially consequential when a retry changes the response condition~\cite{kapoor2026hal,zhang2026dece}. An attempt is one recorded response within a decision. The engine applies one resolved action for that decision; several decisions can occur during one player's turn.

The finite cohort is our population of description. Seat balance does not make the games independent draws from a common distribution of opportunities. Decisions within a game share an evolving board, and rotations within a seed block share initial random sequences. Purchases, jail choices and eliminations alter how those sequences are consumed. Two saved treatment fingerprints also divide games 1--4 from games 5--20. The inspected source diff changes campaign-concurrency plumbing and other application files without changing engine, prompt, resolver or model-client files. This inspection does not identify an effect, or establish the absence of one. We therefore report observed outcomes and paths without treating thousands of decisions as independent evidence of an intrinsic model ranking.

Our analysis combines whole trajectories with narrower opportunity-conditioned summaries. This follows the general motivation for process-sensitive agent evaluation~\cite{ma2024agentboard}, but requires explicit economic units. A deed acquisition changes ownership of one property. A group completion occurs when a player acquires the last deed needed for a color group. A trade is an atomic settlement between two players, with its transfers and charges evaluated together. We distinguish the original thread initiator, the author of the final offer and its accepting player. These roles need not coincide. We record completion at stable transaction boundaries so that an intermediate transfer inside a swap cannot create a fictitious opportunity.

\subsection{Opportunities, expenditures and observed receipts}

For each trade-completed group, the first management window begins at the owner's first actual optional-management decision after settlement and includes further management decisions in that same global turn. We distinguish construction legal on entry, construction becoming legal later in that window, and construction actually applied. Jail and other intervening phases can postpone management. Financing choices within the window can create a build opportunity. Thus, conditioning on entry buildability answers a different question from following every acquisition. It does not remove the consequences of earlier play.

We measure negotiation initiation using two denominators. For configured agent $m$, let $\mathcal{T}_m$ contain its own turns with at least one menu permitting a proposal, and let $\mathcal{P}_m\subseteq\mathcal{T}_m$ contain turns with an applied proposal. If $N_m$ is its number of proposals, then
\begin{equation}
 q_m=\frac{|\mathcal{P}_m|}{|\mathcal{T}_m|},
 \qquad
 i_m=\frac{N_m}{|\mathcal{P}_m|}.
 \label{eq:sol-proposals}
\end{equation}
The first quantity describes proposing-turn propensity; the second describes proposals per proposing turn. Neither is bargaining quality. A denominator of eligible decisions would also count repeated management requests caused by earlier mortgages or purchases, rather than additional turns. We retain that count for auditability without equating it with independent access.

A spending unit $u=(g,p,c)$ joins game $g$, player $p$ and color group $c$ with positive own expenditure on houses or hotels. Including hotels prevents the omission of an owner who inherits four-house deeds and purchases only hotel conversions. Let $\tau_u$ be the first own construction event and $T_g$ the game's terminal event. We report
\begin{equation}
 K_u=\sum_{e\in\mathcal{B}_u}\operatorname{cost}(e),
 \qquad
 R_u^{\mathrm{post}}=
 \sum_{\substack{e\in\mathcal{R}_u\\\tau_u\leq\operatorname{seq}(e)\leq T_g}}
 \operatorname{paid}(e).
 \label{eq:sol-receipts}
\end{equation}
Here $\mathcal{B}_u$ contains that owner's construction events and $\mathcal{R}_u$ contains paid rent on that owner's developed deeds in the group. We also report whole-game developed receipts, without the $\tau_u$ filter. All amounts are game dollars. Acquisition costs, redemption charges, building salvage, mortgage proceeds and estate transfers remain separate flows. An unpaid terminal rent demand is not a paid receipt.

These quantities do not estimate incremental investment returns. The owner can inherit improvements, and subsequent rent depends on opponents' movement and survival. Removing receipts before first own spending fixes a timing mismatch while leaving inherited productive assets in the later window. Observation ends at ownership loss or game end. A late group and an early group consequently have different exposure. We preserve the sequence of acquisitions and expenditures instead of treating $R_u^{\mathrm{post}}/K_u$ as a comparable return on investment. Similar Monopoly investment questions appear in FinBoardBench; our event accounting addresses realized flows and legal financing rather than a universal valuation of the acquired assets~\cite{hu2026finboardbench}.

Ownership, physical development, cash and bank stock are reconstructed from events. The retained audit reconciles these quantities with all 6,989 initial decision payloads and checks construction and salvage against cash-change events. We distinguish a group's completion route from the acquisition routes of its constituent deeds. A group can contain inherited deeds but become complete through a later purchase or trade. No spending unit in this cohort has a second complete-ownership interval at the observed decision or terminal boundaries.

\subsection{Local legal-sequence tests}

To test a financing or construction claim, we restore the actual historical decision prefix, check its event agreement, and execute a specified sequence of the acting player's own legal actions through the engine. We stop at a declared local endpoint, such as completing a proposed construction multiset or paying the current debt. A test cannot borrow another player's management opportunity. Conditional acceptance, altered stock or other diagnostic changes must be stated separately. Legal-action interfaces and replay already occur in related economic benchmarks~\cite{wu2026agentmonopoly,zhang2026dece}; the present use is to adjudicate concrete claims about recorded trajectories.

For a debtor with cash $C$, a generous nominal upper bound is
\begin{equation}
 U=C+V_{\mathrm{all\ building\ salvage}}
       +V_{\mathrm{all\ clear\ deed\ mortgages}}.
 \label{eq:sol-upper}
\end{equation}
The mortgage term includes collateral that clearing buildings could unlock. If $U$ is below debt $D$, own-asset liquidation cannot pay $D$. If $U\geq D$, a legal sequence must still establish access. A hotel downgrade needs bank houses, and the even-selling rule can prevent reaching nominal salvage. Conversely, a current menu can understate finance because mortgages become available only after sales. A tested payment route establishes current-debt feasibility, not future survival or the value of surrender. Once a branch changes the state, the historical policy suffix does not identify how any player would continue.

\section{Economic behavior across complete games}
\label{sec:sol-economic-results}

\subsection{Acquisition and bargaining create several routes to victory}

Across the cohort, the event accounting records 475 direct, 72 auction, 207 trade and 397 inherited deed acquisitions. These are ownership transitions, including repeated acquisitions of the same deed. The distinct group-completion denominator contains 209 episodes: sixteen direct-purchase completions, nine auction completions, sixty-three trade completions and 121 inheritance completions. Sixty-nine inherited completions occur at game-ending bankruptcy and offer no subsequent opportunity. Counting all complete groups as productive achievements would therefore overstate what the surviving agent had occasion to use.

Negotiation varies substantially in activity and role. Sol proposes on 83 of 834 eligible own turns, Gemini on sixteen of 766, Opus on 160 of 660 and Grok on 106 of 676. Their proposal totals are 132, eighteen, 240 and 136, respectively. Gemini's low initiation coexists with twenty-six applied acceptances. These acceptances count the accepting actor, rather than both sides of a settlement. Recipient bargaining can therefore matter even when a participant seldom opens a thread. Conversely, repeated proposals can reflect an extended search for an acceptable exchange without producing a settlement.

The 113 settlements complete a color group for neither party in sixty cases, one party in forty-three, and both parties in ten. In forty-four settlements, the final offer's author differs from the original thread initiator. Bilateral reconstruction is necessary both for interpreting bargaining and for avoiding an attribution error in which the initiating agent receives credit for the other party's counteroffer. Group completion is one observable consequence. A transaction can also exchange railroad income, clear or mortgaged deeds, cash reserves and blocking positions. It can finance the seller's construction while giving the buyer a complete group. A completion count consequently cannot decide which side benefits, much less supply a welfare score.

The G20 New York negotiation illustrates the distinction between a chosen budget and absolute infeasibility. At d235, accepting the actual \$400 offer would leave \$97 after transfer interest. Tested spare mortgages, redemption and four orange houses leave \$41 while preserving brown hotels. Six houses are also reachable with \$91 left if the player liquidates brown development. Later \$250 terms permit six houses with \$64 remaining while preserving the browns. The earlier assertion of no reachable development is too strong, but retaining an existing income source is a real allocation consideration. The branch does not establish that accepting the earlier terms was preferable. It clarifies the assets and income sacrificed by each feasible budget.

Two winners show why the economic account cannot be organized around successful bargaining alone. G14 Sol and G15 Gemini participate in no completed trade. G14 Sol acquires the remaining orange by purchase, constructs six orange houses, mortgages States Avenue, St. Charles Place and Water Works for \$215, and adds three further houses in the same turn. Its later portfolio expands through creditor inheritance. In G15, Gemini first completes and develops oranges in turn 47, following its purchase of Tennessee. Oranges collect \$3,250 in developed rent and trigger the first two eliminations. Inherited browns later trigger the final elimination; inherited dark blues collect no rent and trigger none. Greens arrive only in the terminal estate, with no subsequent action opportunity. Gemini makes no proposals in fifty-three eligible management decisions. These trajectories demonstrate successful acquisition and deployment without completed trades in the observed games. They do not establish the expected performance of abstaining from trade, since purchase access, opponents and subsequent movement differ.

\subsection{Development reflects financing, reserve choices and execution}

All sixty-three trade-completed group episodes have an observed first management window. Forty-three receive construction in that window. Forty-two are buildable on entry, and thirty-four of those are developed during the window. Fifty-four become buildable at some point within it. The difference between entry and later legality matters because a player may first mortgage another deed or redeem the acquired property. Neither the 43/63 whole-path summary nor the 34/42 conditional summary measures willingness or general conversion skill. The first retains the consequences of acquisition and financing; the second describes responsiveness within an already available opportunity.

The eight entry-buildable groups without construction have different recorded explanations. Three end through engine defaults, three through received END actions that conflict with accompanying content, and two through consistent END records with stated reserve rationales. These are eight group episodes, not eight independent experiments. A conflict is not automatically a discarded executable build. One accompanying allocation is uneven as encoded, and another asks for finance. The other twenty-one groups are not entry-buildable, so they cannot simply be added to a tally of declined investments. Following the actual window separates groups that gain access from those that never do.

The construction audit explains one repeated execution failure in greater detail. An even final allocation can be illegal when submitted as consecutive grouped purchases, because intermediate states must also satisfy even building. All sixty-one Opus defaulted purchase multisets execute after interleaving the identical purchases. Sixty chains contain a refusal followed by an invalid action-only retry; G2 d129 has the invalid attempt before the refusal. The exhausted decision then applies END. Some plans contain hotels and require finishing construction in one group to return houses for another. The defect is execution order rather than the affordability or final allocation of those tested plans.

This mechanism is relevant to capability as deployed. The configured participant fails to realize an affordable allocation that it has represented in an attempted action. It does not identify how a different retry policy would change winning chances. Other participants also illustrate why recovery needs a finer definition than eventual legal construction. Among successful decisions following an uneven-building rejection, the original purchase multiset is preserved in 20/25 Sol, 6/18 Gemini and 5/5 Grok recoveries. The remaining successful recoveries change it. G2 d243's original seven-house target is legal after interleaving and leaves \$208; the actual retry and next decision buy six houses and leave \$358. In contrast, G14 d145--148 reaches the original eight-house target within the same global turn. Within-decision preservation and same-turn target attainment answer different questions.

Chosen restraint remains a substantive economic behavior. At G1 d426, Grok holds \$488 with one house on each yellow and applies a consistent END. Its corrected financing account reaches \$1,418. A legal extra Atlantic house would cost \$150, and subsequent liquidation would reach \$1,343. The \$75 difference is the purchase's resale loss. The hold preserves a measurable financing buffer without proving that it maximizes expected value. At G3 d176, Opus's valid action-only retry explicitly favors keeping \$204 over one \$150 house and a \$54 reserve; the record marks no fallback. A prior refusal therefore cannot, by itself, turn every later hold into an engine-authored default.

\subsection{Construction, realized income and inherited portfolios}

The 103 spending units incur \$103,450 in own construction expenditure. Whole-game paid developed rent to those units totals \$93,000, of which \$92,550 follows the owner's first own expenditure. The \$450 difference is G2 light-blue rent collected on inherited development before the winner's first own \$200 construction. Even the later window includes rent supported by inherited buildings. The totals describe expenditure and receipts on observed ownership paths; subtracting them does not yield net profit or construction's causal contribution.

Forty spending units collect no paid developed rent. Twenty-six of those also produce no creditor elimination on that group. These counts concern game/player/color units, whereas fifty-seven of sixty elimination events arise from rent on developed color properties. A large unpayable demand can end an opponent's game and transfer assets without appearing as a paid-rent receipt. Other zero-rent units retain assets, recover salvage, or face little remaining exposure. Denial and deterrence are possible sources of value but are not measured here. The data therefore support zero observed receipts, rather than forty failed investments.

Successive owners of G19's dark blues provide a concrete contrast. Opus spends \$1,200, receives \$175 in developed rent and later realizes \$600 in building sales. Gemini inherits the deeds, spends \$2,000 over its subsequent ownership, and receives \$1,400 in developed rent on a later Sol landing. The two units concern the same properties with different timing, financing and exposure. They show why an attractive group cannot guarantee receipts and why inherited ownership must remain visible. They do not isolate an owner effect. A recent \$1,200 construction step in the second trajectory must not be substituted for Gemini's \$2,000 cumulative own expenditure.

Winners account for fifty-eight spending units. Thirty-four contain at least one inherited deed at first own expenditure, while thirty first become complete through an estate. G3 brown, G6 light blue, G12 light blue and G14 green account for the distinction: inherited pieces are present, but later purchases or trades close the groups. A HOUSE\_BUILT-only definition would also miss G10's winner, which inherits four-house light blues and purchases \$150 in direct hotel conversions. The spending definition must include that case even though it records no own house purchase on the group.

Creditor inheritance is a major feature of the realized concentration of assets. The eventual winner is creditor in fifty-four of sixty eliminations, including thirty-five of forty nonterminal eliminations. Nineteen of twenty final eliminations name the winner as creditor; G13 ends with Luxury Tax owed to the bank. The nonterminal count is useful because it excludes estates received only when play has already ended. Under this implementation's declared variant, developed property can pass intact to a creditor. Productive assets accumulated by one player can consequently support another's later income and construction. The observed concentration is jointly determined by prior play, debts, movement and this rule. It is not a causal estimate of inheritance's advantage.

\subsection{Debt payment requires accessible finance}

The elimination denominator comprises all sixty applied bankruptcies. Fifty-six follow declarations, and four occur automatically after mortgages exhaust the permitted continuation without covering the debt. In fifty-six cases, even the generous upper bound in Equation~\ref{eq:sol-upper} is below the current debt. Three current debts have engine-verified payment routes. The remaining case has enough nominal value but a binding hotel-liquidation restriction. Table~\ref{tab:sol-financing} states the local endpoints rather than inferring later survival.

\begin{table}[t]
\centering
\small
\begin{tabular}{llrrp{0.34\linewidth}}
\toprule
Case & Agent & Debt & Nominal upper & Local legal result \\
\midrule
G12 d190 & Sol & \$500 & \$748 & Payment route retains \$8. \\
G12 d249 & Grok & \$250 & \$613 & Five house sales pay; \$43 remains. \\
G17 d234 & Gemini & \$550 & \$905 & Sales and two mortgages pay; \$155 remains. \\
G7 d339 & Opus & \$750 & \$1,310 & Only \$440 is accessible; hotel sale is locked. \\
\bottomrule
\end{tabular}
\caption{Four eliminations not ruled out by the nominal upper bound. Payment routes need not realize the entire upper bound. The other fifty-six eliminations cannot cover the current debt even under that bound. All sixty cases exclude assumed assistance from another player.}
\label{tab:sol-financing}
\end{table}

The payable cases do not share a single reasoning mechanism. At G12 d190, Sol's stated \$408 maximum omits red collateral unlocked by selling houses. At G17 d234, Gemini's \$530 account understates a route involving sales and Baltic and Park Place mortgages. G12 d249 instead contains a sharp action/explanation disagreement. Its accompanying content requests house sales, and its final provider and applied private explanations recommend survival, while the received tool declares bankruptcy. Five legal one-house sales already cover the debt without a mortgage. The applied elimination is avoidable at that local endpoint, but the record does not identify an intention to surrender or justify calling all three cases deliberate choices.

Successful liquidation supplies a necessary contrast. G18 d428 sells development and mortgages red collateral, pays \$750 and retains \$69. G20 d327--336 starts with \$49, realizes \$700 in salvage and \$665 in mortgages, pays \$1,400, and retains \$14 using the actual stock constraints. These paths exhibit staged financing and preservation of continued participation. At G16 d602, a twenty-three-step legal liquidation reaches \$1,393 against \$2,000 owed, so even extensive correct execution cannot pay. The applied private note already reports the correct total; a generated research label instead gives \$1,233 by omitting \$160 of unlocked light-blue collateral. Legal-sequence checking corrects the analyst as well as the agent.

Finite stock creates a different limit. Nine games reach zero bank houses in initial decision payloads, but an empty bank does not establish an agent's desire for a blocked project. At G7 d339, Opus has pinks at four houses, four houses and a hotel. Too few bank houses prevent downgrading the hotel, and even-selling prevents selling a lower sibling first. Its \$390 cash and Oriental mortgage yield only \$440 against \$750, despite \$1,310 nominal value. Exhausting both permitted mortgage orders leaves the sale menu empty. This establishes a genuine current-debt lock. Earlier opportunities to convert hotels and release houses belong to different states and cannot be carried forward as an escape from that debt.

Other scarcity cases allow costly adjustment. At G12 d312, a direct Illinois hotel is blocked with reds at 3/3/4 and no bank houses, yet a longer legal restructuring reaches it with \$26 and four bank houses remaining. That route changes other development and reserves. At G5 d300, a legal \$100 pink hotel would release four houses, but retaining the houses does not by itself establish that a funded rival project was denied. These cases distinguish immediate restrictions, reachable reallocation and expressed demand. Economic planning in this cohort consists of choosing among such consequential allocations, sometimes executing them successfully, sometimes preserving reserves, and sometimes losing an accessible action through the response process. The accompanying bargaining statements require a separate reconstruction of what counterparties saw and how they responded.
